\section{Key Architecture}

\label{sec:keys}

\subsection{Root Key Derivation}

The user's root key is derived from a passphrase using Argon2id
\cite{argon2}:

\begin{equation}
  \text{rootKey} = \text{Argon2id}(\text{passphrase}, \text{salt}, t=3, m=65536, p=4)
\end{equation}

Parameters: 3 iterations, 64 MiB memory, 4 parallel lanes.  The salt is
generated once and stored alongside the encrypted vault header (the salt
is not secret).  The output is 32 bytes, used as the master symmetric key.

The root key never leaves the device.  It is held in memory only during
an active session and zeroed on lock/logout.

\subsection{Per-Vault Data Encryption Key}

Each vault (a logical container for related data---e.g., a company's
cap table, a personal document set) has its own DEK.  The DEK is a random
256-bit key, encrypted to the vault owner's ML-KEM-1024 public key:

\begin{equation}
  \text{encDEK} = \text{ML-KEM-1024.Encapsulate}(\text{ownerPubKey}) \to (\text{encapsulated}, \text{sharedSecret})
\end{equation}

The shared secret serves as the DEK.  The encapsulated key blob is stored
alongside the vault metadata.

To share a vault with another user, the DEK is re-encapsulated to the
recipient's ML-KEM-1024 public key.  Each recipient gets their own
encapsulated key blob; the DEK itself is never transmitted in plaintext.

\subsection{Per-Document Encryption}

Individual documents are encrypted with AES-256-GCM using the vault DEK
and a random 12-byte nonce per document.  This is implemented in
\texttt{lux/transfer/pkg/crypto/envelope.go}:

\begin{lstlisting}[language=Go,caption={envelope.go: ML-KEM + AES-256-GCM envelope}]
func Seal(plaintext []byte,
          recipientPubKey []byte) (*Envelope, error) {
    pk, err := mlkem.PublicKeyFromBytes(
        recipientPubKey, mlkem.MLKEM1024)
    // ML-KEM encapsulate: shared secret + encapsulated key
    encapsulated, sharedSecret, err := pk.Encapsulate()
    // AES-256-GCM with shared secret as key
    block, err := aes.NewCipher(sharedSecret[:32])
    gcm, err := cipher.NewGCM(block)
    nonce := make([]byte, gcm.NonceSize())
    rand.Read(nonce)
    ciphertext := gcm.Seal(nil, nonce, plaintext, nil)
    return &Envelope{
        Ciphertext:      ciphertext,
        EncapsulatedKey: encapsulated,
        Nonce:           nonce,
        Algorithm:       AlgoMLKEM1024_AES256GCM,
    }, nil
}
\end{lstlisting}

Decryption requires the recipient's ML-KEM-1024 private key to
decapsulate the shared secret, then AES-256-GCM decryption:

\begin{lstlisting}[language=Go,caption={envelope.go: decryption path}]
func Open(env *Envelope,
          recipientPrivKey []byte) ([]byte, error) {
    sk, err := mlkem.PrivateKeyFromBytes(
        recipientPrivKey, mlkem.MLKEM1024)
    sharedSecret, err := sk.Decapsulate(env.EncapsulatedKey)
    block, err := aes.NewCipher(sharedSecret[:32])
    gcm, err := cipher.NewGCM(block)
    plaintext, err := gcm.Open(
        nil, env.Nonce, env.Ciphertext, nil)
    return plaintext, nil
}
\end{lstlisting}

\subsection{Post-Quantum Signatures}

All envelopes can be signed using ML-DSA-87 (FIPS 204 \cite{fips204}).
The signer's public key is included in the envelope for verification:

\begin{lstlisting}[language=Go,caption={envelope.go: ML-DSA-87 signing}]
func Sign(data []byte,
          signerPrivKey []byte) ([]byte, error) {
    sk, err := mldsa.PrivateKeyFromBytes(
        mldsa.MLDSA87, signerPrivKey)
    return sk.Sign(data)
}
\end{lstlisting}

ML-DSA-87 provides NIST Security Level 5---equivalent security to
AES-256 against quantum adversaries.

\subsection{Algorithm Summary}

\begin{table}[h]
\centering
\small
\begin{tabular}{@{}llll@{}}
\toprule
\textbf{Function} & \textbf{Algorithm} & \textbf{Standard} & \textbf{Security Level} \\
\midrule
Key derivation    & Argon2id       & RFC 9106       & Memory-hard \\
Key encapsulation & ML-KEM-1024    & FIPS 203       & NIST Level 5 \\
Symmetric encrypt & AES-256-GCM    & FIPS 197       & 256-bit \\
Digital signature  & ML-DSA-87      & FIPS 204       & NIST Level 5 \\
Threshold ECDSA   & CGGMP21        & ---            & secp256k1 \\
Threshold EdDSA   & FROST          & RFC 9591       & Ed25519 \\
\bottomrule
\end{tabular}
\caption{Cryptographic algorithms and their security levels.}
\label{tab:algorithms}
\end{table}

% ======================================================================
